\chapter{Python-Schnittstelle}\label{ch:python}

Das Paket \texttt{fouriergrating} kann alles, was \cmd{fgsolve} kann, ohne Umweg über Dateien. Die Konfiguration ist
dasselbe JSON-Schema als Python-\texttt{dict}; Einheiten wie in JSON (\textmu m, Grad). Die Bindungen enthalten keine
eigene Physik: Ergebnisse sind identisch mit \cmd{fgsolve}.

\section{Überblick}
\begin{center}
\small
\begin{tabularx}{\linewidth}{lY}
\toprule
Funktion / Klasse & Zweck \\
\midrule
\texttt{fg.load\_config(pfad)} & JSON-Datei lesen; \key{output.dir} und \key{cad.file} relativ zur Datei auflösen \\
\texttt{fg.solve(cfg, profile=False)} & erster Scanpunkt $\to$ \texttt{Result} (\key{sweep} wird ignoriert) \\
\texttt{fg.solve\_scan(cfg, profile=False)} & alle Sweep-Werte und Scanpunkte (parallel) $\to$ \texttt{ScanResult} \\
\texttt{fg.near\_field(cfg)} & Nahfeldschnitte des ersten Scanpunkts \\
\texttt{fg.Simulation(cfg)} & vorbereitete Rechnung für viele Auswertungen, Parameteränderungen, Gradienten \\
\texttt{fg.set\_backend(pref, threshold)} & Standard-Backend für Konfigurationen ohne \key{backend} \\
\texttt{fg.get\_backend()}, \texttt{fg.cuda\_available()}, \texttt{fg.core\_variant()} & Abfragen \\
\texttt{fg.cad.check}, \texttt{fg.cad.preview}, \texttt{fg.cad.llm\_tool\_spec} & CAD (Kapitel~\ref{ch:cad}) \\
\texttt{fouriergrating.torch}, \texttt{fouriergrating.jax} & Gradienten-Adapter (Kapitel~\ref{ch:optimierung}) \\
\texttt{fg.FourierGratingError} & Ausnahme für Konfigurations- und Rechenfehler (Unterklasse von \texttt{RuntimeError}) \\
\bottomrule
\end{tabularx}
\end{center}

\section{Ergebnisse: \texttt{Result}}
\begin{lstlisting}[style=py]
res = fg.solve(cfg)
res.reflected        # numpy structured array, Felder: n, m, eff, eff_s, eff_p, amp_s, amp_p
res.transmitted
res.efficiency("T", -1)        # Effizienz einer Ordnung (0, falls nicht propagierend)
res.efficiency("R", 0, 1)      # 2D: Ordnung (0, 1)
res.absorption
res.diagnostics      # energy_balance, max_condition_t11, rk_steps, warnings, precision, timings, ...
res.incidence        # {"lambda_um", "phi_deg", "psi_deg", "delta_deg"}
res.scattering       # isolierte Strukturen: {"up", "down", "lateral", "absorbed"}
res.modal            # Modenanregung: {"incident", "reflected", "transmitted", "radiated"}
res.cylindrical      # Zylinder: {"reflected", "transmitted", "absorbed", "harmonics", ...}
\end{lstlisting}
Die Felder \texttt{amp\_s}, \texttt{amp\_p} sind komplex. Beispiel: alle reflektierten Ordnungen ausgeben
\begin{lstlisting}[style=py]
for o in res.reflected:
    print(f"R({o['n']:+d},{o['m']:+d})  e = {o['eff']:.5f}  s = {o['eff_s']:.5f}  p = {o['eff_p']:.5f}")
\end{lstlisting}

\section{Scans: \texttt{ScanResult}}
\begin{lstlisting}[style=py]
scan = fg.solve_scan(cfg)
scan.points           # Liste von Result (mit .incidence)
scan.sweep_parameter  # Pointer des Sweeps oder ""
scan.columns          # dict von numpy-Arrays wie efficiencies.csv:
                      # sweep_value, lambda_um, phi_deg, psi_deg, delta_deg, side, n, m,
                      # efficiency, eff_s, eff_p, amp_s, amp_p
df = scan.to_pandas() # benötigt pandas
\end{lstlisting}
Spektrum der Ordnung $T_{-1}$ zeichnen:
\begin{lstlisting}[style=py]
import matplotlib.pyplot as plt
c = scan.columns
sel = (c["side"] == "T") & (c["n"] == -1)
plt.plot(c["lambda_um"][sel], c["efficiency"][sel])
\end{lstlisting}

\section{\texttt{Simulation}: wiederholte Rechnungen}\label{sec:simulation}
\texttt{fg.Simulation} baut die Schichtmodelle (Fourierkoeffizienten, Normalenfelder) einmal auf und behält sie.
Danach sind Rechnungen mit anderem Einfall deutlich schneller, und einzelne Parameter lassen sich ändern.
\begin{lstlisting}[style=py]
sim = fg.Simulation(cfg)
r1 = sim.solve(lambda_um=0.6328, phi_deg=30.0)   # fehlende Werte: erster Scanpunkt
r2 = sim.solve(lambda_um=0.55)
sim.set_parameter("/layers/0/thickness_um", 0.25) # JSON-Pointer wie beim Sweep
r3 = sim.solve()
sim.model_builds      # wie oft Schichtmodelle gebaut wurden
sim.config            # aktuelle Konfiguration als dict
\end{lstlisting}
\begin{itemize}
  \item \texttt{set\_parameter} baut nur die geänderten Schichten neu, wenn sich außerhalb von \key{layers} nichts
        ändert; sonst alle. Das Ergebnis ist bitgleich zu einer neuen \texttt{Simulation}.
  \item Ein \texttt{Simulation}-Objekt ist \textbf{nicht threadsicher}. Mehrere Python-Threads dürfen parallel
        rechnen, jeder mit eigenem Objekt; während der Rechnung ist der GIL freigegeben.
\end{itemize}

\section{Schleifen in Python}\label{sec:py-schleifen}
\subsection{Zweidimensionale Parameterstudie}
\begin{lstlisting}[style=py]
import numpy as np, fouriergrating as fg
cfg = fg.load_config("examples/v3_lamellar_dielectric.json")
sim = fg.Simulation(cfg)
depths = np.linspace(0.1, 0.8, 15)
widths = np.linspace(0.2, 0.8, 13)
E = np.zeros((len(depths), len(widths)))
for i, h in enumerate(depths):
    sim.set_parameter("/layers/0/thickness_um", h)
    for j, w in enumerate(widths):
        sim.set_parameter("/layers/0/breakpoints_um/1", w)
        E[i, j] = sim.solve().efficiency("T", -1)
\end{lstlisting}

\subsection{Littrow-Kurve}
\begin{lstlisting}[style=py]
d = cfg["period_um"]["x"]
for lam in np.linspace(0.5, 0.9, 41):
    phi = np.degrees(np.arcsin(lam / (2 * d)))
    print(lam, sim.solve(lambda_um=lam, phi_deg=phi).efficiency("R", -1))
\end{lstlisting}

\subsection{Dispersive Materialien von Hand}\label{sec:py-dispersion}
Solange \FG{} keine wellenlängenabhängigen Materialien kennt, setzen Sie die Brechzahl je Wellenlänge selbst. Die
Brechzahl muss dafür als \val{[Re, Im]} in der Konfiguration stehen:
\begin{lstlisting}[style=py]
import numpy as np
# Tabelle (lambda in µm, n, k), z. B. aus einer CSV-Datei
tab = np.loadtxt("au_nk.csv", delimiter=",", skiprows=1)
n_of = lambda lam: np.interp(lam, tab[:, 0], tab[:, 1])
k_of = lambda lam: np.interp(lam, tab[:, 0], tab[:, 2])

cfg["materials"]["au"] = {"n": [1.0, 1.0]}       # Platzhalter, wird überschrieben
sim = fg.Simulation(cfg)
for lam in np.linspace(0.5, 0.9, 81):
    sim.set_parameter("/materials/au/n/0", n_of(lam))
    sim.set_parameter("/materials/au/n/1", k_of(lam))
    r = sim.solve(lambda_um=lam)
    print(lam, r.efficiency("R", -1))
\end{lstlisting}
Für eine Sellmeier-Formel ersetzen Sie die Interpolation durch die Formel. Eine Materialänderung baut alle Schichten
neu; die Rechnung ist dadurch nicht teurer als einzelne Aufrufe von \texttt{fg.solve}.

\section{S-Matrix in der s/p-Basis}\label{sec:smatrix}
\begin{lstlisting}[style=py]
S = sim.smatrix(lambda_um=0.6, phi_deg=17, psi_deg=10)   # ein Eigenproblem je Richtung
S.orders, S.alpha, S.gamma, S.beta_sup, S.beta_sub         # Ordnungen und Wellenzahlen
r = S.index(-1, 0)                                          # Index der Ordnung (n, m)
nf = len(S.orders)
a_in = np.zeros(2 * nf, complex); a_in[S.index(0, 0)] = 1.0  # s-polarisiert, Ordnung (0, 0)
refl = S.R_top @ a_in                                       # s bei r, p bei nf + r
\end{lstlisting}
Die vier Blöcke \texttt{R\_top} (von oben zurück nach oben), \texttt{T\_down} (von oben ins Substrat),
\texttt{T\_up} (von unten ins Superstrat) und \texttt{R\_bottom} beschreiben beide Polarisationen und Licht von
oben und unten in einem Aufruf. Die Spalte der Ordnung $(0,0)$ ergibt dieselben Amplituden wie \texttt{solve}.
Andere Spalten beschreiben Einfall aus anderen Richtungen der Bloch-Familie, mit einem um die Einfallsordnung
zentrierten Ordnungsfenster (Abweichung im Rahmen des Trunkierungsfehlers). Voraussetzungen: periodisch,
homogenes isotropes Superstrat und Substrat, keine Koordinatentransformation.

\section{Backend und Kern}
\begin{lstlisting}[style=py]
fg.set_backend("auto", threshold=160)    # Standard für Konfigurationen ohne "backend"
fg.get_backend()
fg.core_variant()                        # "cpu" oder "cuda12"
\end{lstlisting}
\texttt{FG\_CORE=cpu} erzwingt den CPU-Kern beim Import.

\section{Fehlerbehandlung}
\begin{lstlisting}[style=py]
try:
    res = fg.solve(cfg)
except fg.FourierGratingError as e:      # Konfigurations- oder Rechenfehler
    print("Fehler:", e)
except NotImplementedError as e:         # spezifiziert, aber noch nicht umgesetzt
    print("Noch nicht verfügbar:", e)
\end{lstlisting}

\section{Tests}
\begin{lstlisting}[style=shell]
pip install .[test]
pytest tests/python
FG_FGSOLVE=build/cpu-release/apps/fgsolve pytest tests/python   # zusätzlich Nahfeld = fgsolve
\end{lstlisting}
